\begin{table}[!h]
\centering
\renewcommand{\arraystretch}{1.15}
\setlength{\tabcolsep}{6pt}
\captionsetup{justification=centering}
\caption{Heterogeneous Effects on Migration Intention}
\begin{tabular*}{\textwidth}{l@{\extracolsep{\fill}}cccccc}
\hline
& \multicolumn{3}{c}{Pop.\ density} & \multicolumn{3}{c}{Female pop.} \\
\cmidrule(l){2-4} \cmidrule(l){5-7}
& T1 & T2 & T3 & T1 & T2 & T3 \\
\hline
Spanish share 1936-1955$\times$Post & 0.027 & -0.031 & 0.043* & 0.035** & -0.072 & -0.076 \\
 & (0.045) & (0.053) & (0.024) & (0.014) & (0.047) & (0.050) \\
Spanish share 1956-1978$\times$Post & -0.060 & -0.141 & 0.032 & 0.145** & 0.244 & 0.095 \\
  & (0.128) & (0.273) & (0.105) & (0.062) & (0.165) & (0.143) \\
\addlinespace
Observations & 2,556 & 1,936 & 3,103 & 2,841 & 1,753 & 2,635 \\
$p$-value ($\beta_{36{-}55} = \beta_{56{-}78}$) & 0.605 & 0.696 & 0.926 & 0.142 & 0.126 & 0.373 \\
Tercile upper bound & $0.493$ & $0.097$ & $0.546$ & $0.158$ & $0.685$ & $0.278$ \\
\hline
\end{tabular*}
\addvspace{0.2em}
\captionsetup{font=footnotesize, justification=justified, singlelinecheck=false}
\caption*{\footnotesize Notes: The dependent variable is migration intention. Each column reports subsample estimates for the tercile of the specified municipal-level variable. Terciles are defined across municipalities using the pre-treatment average of each municipality-level aggregate over the 2012, 2014, 2017, and 2019 survey waves. Tercile upper bounds are shown at the bottom of each panel. The tercile variables are defined as follows: the share of respondents who are very interested in politics and the share of respondents who have a partner, both calculated at the municipality level by pooling individual-level LAPOP responses across all pre-reform survey rounds. All specifications include municipality and year fixed effects and individual controls for age and a male indicator. The reported coefficients are average marginal effects from a logit specification, expressed as the change in the probability of reporting migration intention per one-unit increase in the corresponding regressor, where one unit corresponds to one percentage point of the municipal Spanish-born population share. Standard errors are clustered at the municipality level and reported in parentheses (56 clusters). The row labeled $p$-value: $\beta^{1936-1955}=\beta^{1956-1978}$ reports the p-value from a two-sided Wald test of equality between the two cohort-specific coefficients within each tercile. Joint Wald tests of the null hypothesis that the coefficient on Spanish share $\times$ Post is equal across the three terciles ($H_0: \beta_{T_1} = \beta_{T_2} = \beta_{T_3}$), computed under the independence of the tercile subsamples and following a $\chi^2$ distribution with 2 degrees of freedom, yield the following p-values, reported as ($\beta^{1936-1955}$ / $\beta^{1956-1978}$) for each variable: Share partnered (0.089 / 0.341); Share very interested in politics (0.060 / 0.578). * $p<0.10$, ** $p<0.05$, *** $p<0.01$.}
\end{table}
